The paper's strongest conclusions are mechanism-level, not universal algorithmic claims.
We therefore state the main boundaries explicitly.

\paragraph{Conditional, not universal, scalar insufficiency.}
The benchmark clearly shows that ordered scalar thresholding is not a sufficient description of prefix-level compute control.
But the downstream ranking among stronger controllers is conditional on domain, scale, and deployment noise.
We do not claim that every one-dimensional encoding is impossible or that scalar rules are always dominated.

\paragraph{Exact-state value does not imply deployable success.}
Exact-state factorization is often strong and sometimes best.
That is useful because it shows the structure is not empty.
But the deployable predicted-state gap remains large, and our evidence indicates that this gap is driven mainly by state identification under noisy deployment.
Targeted rescue probes do not erase this boundary.
Appendix experiments with high-determinacy filtering, pairwise soft targets, exact-state teacher distillation, and their combinations produce partial, domain-specific gains, but they do not surpass the strongest learned one-dimensional baseline on either exact-checker domain.
This is why the paper stops at a benchmark-and-diagnosis claim rather than a broader algorithmic victory claim.

\paragraph{Continuation outcome statistics are not the original \(S_t\).}
In the current exact-checker domains, the original success-uncertainty view of \(S_t\) often degenerates.
The richer continuation features used in the structured controller are therefore better described as continuation outcome statistics.
This preserves the useful intuition of the proposal without overstating what the present experiments establish.

\paragraph{Cross-domain routing is an extension.}
The project also contains a substantial line of cross-domain routing experiments.
Those results are informative, especially as an extension about learned routing versus hard specialist portfolios.
But they answer a different question from the main paper.
We therefore treat them as appendix-level evidence rather than as the central contribution.

\paragraph{Repeatability and scale.}
We include a minimal within-domain repeatability check for the main controller set.
That repeatability result is strong enough to constrain the paper's wording, but it is not the last word on every controller family.
For the main paper, the purpose of repeatability is not to certify a universal winner; it is to prevent us from over-writing single-run frontier effects as stable conclusions.
